\section{Urysohn 引理}\label{section:ch04-urysohn-lemma}

现在我们遇到了本书中的第一个深刻的定理，这个定理通常称为``Urysohn 引理''。它断言：正规空间 $X$ 上存在着某种实值连续函数。这个定理是证明许多重要定理的一个至关重要的工具。其中的三个，即 Urysohn 度量化定理，Tietze 扩张定理，以及关于流形的一个嵌入定理，我们将在本章以后几节中给出证明。

为什么说 Urysohn 引理是一个``深刻的''定理呢？因为它的证明包含着以前的证明中所没有的那种新颖的思想。或许我们可以用以下方式来解释清楚我们的意思：如果将本书所给出的证明通篇删去，然后再把这本书交给一个没有学过拓扑学但很聪明的学生，大体上可以设想，这个学生应该能够通读这本书并且由他自己来完成证明（当然要花费很多时间和精力了，并且不能期望他解决那些棘手的例子）。然而 Urysohn 引理则完全不同，如果不给出详尽的提示，要想掌握这个引理的证明势必要有相当超常的创造性。

\begin{theorem}[Urysohn 引理]\label{theorem:ch04-urysohn-lemma}
设 $X$ 为正规空间，$A$ 和 $B$ 是 $X$ 中两个无交的闭集。$[a,b]$ 是实直线上的一个闭区间。则存在一个连续映射
\begin{align*}
f\colon X\longrightarrow[a,b]
\end{align*}
使得对于 $A$ 中的每一个 $x$，有 $f(x)=a$，并且对于 $B$ 中每一个 $x$，有 $f(x)=b$。
\end{theorem}

\par
\begin{figure}[htbp]
\centering
\includegraphics[width=0.6\textwidth]{ch04-fig33-1.png}
\caption{Urysohn 引理中的开集 $U_p$}\label{figure:ch04-urysohn-u-p}
\end{figure}

\par
\begin{figure}[htbp]
\centering
\includegraphics[width=0.7\textwidth]{ch04-fig33-2.png}
\caption{$f$ 的连续性}\label{figure:ch04-urysohn-continuity}
\end{figure}

\begin{proof}
我们只要就区间 $[0,1]$ 的情形来讨论就够了。一般情形可以由此推出。证明的第一步是应用正规性构造 $X$ 的开集族 $U_{p}$，其中下标是有理数。然后利用这些集合来确定连续函数 $f$。

第一步。设 $P$ 是区间 $[0,1]$ 中的所有有理数构成的集合，对于每一个 $p\in P$，我们定义 $X$ 中的一个开集 $U_{p}$，使得当 $p<q$ 时，有
\begin{align*}
\overline{U}_{p}\subset U_{q}.
\end{align*}
这样，包含关系便是这些集合 $U_{p}$ 间的一个全序，这个全序与下标在实直线上通常的序关系相同。

由于 $P$ 是可数的，我们可以通过归纳原则（确切地说是应用归纳定义原理）来定义这些集合 $U_{p}$。以某种方式将 $P$ 中元素排列成一个无穷序列；为了方便起见，不妨设 $1$ 和 $0$ 就是序列最前面的两个元素。

现在定义集合 $U_{p}$ 如下：首先，令 $U_{1}=X-B$。其次，因为 $A$ 是包含在开集 $U_{1}$ 中的闭集，根据 $X$ 的正规性，可以选取一个开集 $U_{0}$ 使得
\begin{align*}
A\subset U_{0}\quad\text{和}\quad\overline{U}_{0}\subset U_{1}.
\end{align*}

一般地，令 $P_{n}$ 表示有理数序列中前 $n$ 项所构成的集合，假设对于所有属于 $P_{n}$ 的有理数 $p$，开集 $U_{p}$ 都已经定义好了，并且满足条件
\begin{align}
p<q\Longrightarrow\overline{U}_{p}\subset U_{q}.
\end{align}
设 $r$ 表示有理数序列中第 $n+1$ 项，我们来定义 $U_{r}$。

考虑集合 $P_{n+1}=P_{n}\cup\{r\}$。它是区间 $[0,1]$ 的一个有限子集，并且它有一个由实直线上通常的序关系 $<$ 给出的全序。在一个有限的全序集中，每一个元素（除了最小元和最大元外）有一个直接前元和一个直接后元（见有限全序集的基本性质）。这个全序集 $P_{n+1}$ 的最小元是 $0$，最大元是 $1$，$r$ 既不是 $0$ 也不是 $1$，所以 $r$ 在 $P_{n+1}$ 中有一个直接前元 $p$ 和一个直接后元 $q$。集合 $U_{p}$ 和 $U_{q}$ 已有定义，并且根据归纳假设，有 $\overline{U}_{p}\subset U_{q}$。应用 $X$ 的正规性，我们能找到 $X$ 的一个开集 $U_{r}$，使得
\begin{align*}
\overline{U}_{p}\subset U_{r}\quad\text{和}\quad\overline{U}_{r}\subset U_{q}.
\end{align*}

我们断言：对于 $P_{n+1}$ 中的每一对元素，(*) 成立。若这两个元素都属于 $P_{n}$，则由归纳假定可见 (*) 成立。若它们中有一个是 $r$，另一个是 $P_{n}$ 中的点 $s$，那么在 $s\leqslant p$ 的情况下有
\begin{align*}
\overline{U}_{s}\subset\overline{U}_{p}\subset U_{r},
\end{align*}
在 $s\geqslant q$ 的情况下有
\begin{align*}
\overline{U}_{r}\subset\overline{U}_{q}\subset U_{s}.
\end{align*}
于是，对于 $P_{n+1}$ 中每一对元素，(*) 成立。

根据归纳原则，对于所有的 $p\in P$，$U_{p}$ 已有定义。

为更清楚地说明这个过程，假设我们先用标准方式将 $P$ 中的元素排成无穷序列：
\begin{align*}
P=\Big\{1,0,\frac{1}{2},\frac{1}{3},\frac{2}{3},\frac{1}{4},\frac{3}{4},\frac{1}{5},\frac{2}{5},\frac{3}{5},\cdots\Big\}.
\end{align*}
在定义 $U_{0}$ 和 $U_{1}$ 之后，我们定义 $U_{1/2}$，使得 $\overline{U}_{0}\subset U_{1/2}$，且 $\overline{U}_{1/2}\subset U_{1}$。于是在 $U_{0}$ 与 $U_{1/2}$ 之间有相应的 $U_{1/3}$。在 $U_{1/2}$ 与 $U_{1}$ 之间有相应的 $U_{2/3}$ 等等。\cref{figure:ch04-urysohn-u-p} 里所描绘的就是第八步时的情况，而第九步就是在 $U_{1/3}$ 和 $U_{1/2}$ 之间选取开集 $U_{2/5}$。如此等等。

第二步。现在对于区间 $[0,1]$ 中每一个有理数 $p$，我们已定义了 $U_{p}$。我们将上述定义扩充到 $\mathbb{R}$ 中所有的有理数 $p$ 上，方法如下：
\begin{align*}
U_{p}=\varnothing\quad\text{如果 }p<0,
\end{align*}
\begin{align*}
U_{p}=X\quad\text{如果 }p>1.
\end{align*}
可以验证，对于任意一对有理数 $p$ 和 $q$，
\begin{align*}
p<q\Longrightarrow\overline{U}_{p}\subset U_{q}
\end{align*}
仍然是正确的。

第三步。给定 $X$ 的一个点 $x$，我们定义 $\mathbb{Q}(x)$ 为这样一些有理数 $p$ 的集合：$p$ 所对应的开集 $U_{p}$ 包含 $x$，即
\begin{align*}
\mathbb{Q}(x)=\{\,p\mid x\in U_{p}\,\}.
\end{align*}
因为对于 $p<0$，没有点 $x$ 属于 $U_{p}$，所以 $\mathbb{Q}(x)$ 不包含小于 $0$ 的数。又因为对于 $p>1$，每一个 $x$ 都在 $U_{p}$ 中，所以 $\mathbb{Q}(x)$ 包含大于 $1$ 的每一个数。因此，$\mathbb{Q}(x)$ 是有下界的，并且它的下确界是区间 $[0,1]$ 中的点。于是规定
\begin{align*}
f(x)=\inf\mathbb{Q}(x)=\inf\{\,p\mid x\in U_{p}\,\}.
\end{align*}

第四步。证明 $f$ 就是所求的函数。若 $x\in A$，则对于每一个 $p\geqslant 0$，有 $x\in U_{p}$，于是 $\mathbb{Q}(x)$ 等于非负有理数集，所以 $f(x)=\inf\mathbb{Q}(x)=0$。类似地，若 $x\in B$，因没有 $p\leqslant 1$ 使得 $x\in U_{p}$，于是 $\mathbb{Q}(x)$ 由所有大于 $1$ 的有理数组成。因此 $f(x)=1$。

所有这些都是容易的，唯一困难的部分是证明 $f$ 的连续性。为此，我们先证明以下几个基本事实：
\begin{enumerate}
\item $x\in\overline{U}_{r}\Longrightarrow f(x)\leqslant r$.
\item $x\notin U_{r}\Longrightarrow f(x)\geqslant r$.
\end{enumerate}

为了证明 (1)，注意这样一个事实：如果 $x\in\overline{U}_{r}$，则对任何 $s>r$，$x\in U_{s}$。因此，$\mathbb{Q}(x)$ 包含所有大于 $r$ 的有理数。于是根据定义有
\begin{align*}
f(x)=\inf\mathbb{Q}(x)\leqslant r.
\end{align*}
再来证明 (2)。注意这样一个事实：如果 $x\notin U_{r}$，则对于任何 $s<r$，$x$ 不在 $U_{s}$ 中。因此，$\mathbb{Q}(x)$ 不包含小于 $r$ 的有理数，于是
\begin{align*}
f(x)=\inf\mathbb{Q}(x)\geqslant r.
\end{align*}

现在我们来证明 $f$ 的连续性。给定 $X$ 的一个点 $x_{0}$ 和 $\mathbb{R}$ 中包含点 $f(x_{0})$ 的开区间 $(c,d)$。要找 $x_{0}$ 的一个邻域 $U$，使得 $f(U)\subset(c,d)$。取有理数 $p$ 和 $q$，使得
\begin{align*}
c<p<f(x_{0})<q<d.
\end{align*}
我们断言开集
\begin{align*}
U=U_{q}-\overline{U}_{p}
\end{align*}
就是所要找的点 $x_{0}$ 的邻域。如\cref{figure:ch04-urysohn-continuity}所示。

首先，注意 $x_{0}\in U$。因为根据 (2) 及 $f(x_{0})<q$ 可见 $x_{0}\in U_{q}$。同时，再根据 (1) 及 $f(x_{0})>p$ 可见 $x_{0}\notin\overline{U}_{p}$。

其次，我们证明 $f(U)\subset(c,d)$。令 $x\in U$。这时有 $x\in U_{q}\subset\overline{U}_{q}$，因此根据 (1) 可见 $f(x)\leqslant q$。又由于 $x\notin\overline{U}_{p}$，所以 $x\notin U_{p}$，并且根据 (2) 有 $f(x)\geqslant p$。从而，$f(x)\in[p,q]\subset(c,d)$。定理证毕。
\end{proof}
\begin{definition}\label{definition:ch04-separated-by-function}
设 $A$ 和 $B$ 是拓扑空间 $X$ 的两个子集，如果存在一个连续函数 $f\colon X\to[0,1]$，使得 $f(A)=\{0\}$ 及 $f(B)=\{1\}$，那么就称 $A$ 和 $B$ 能用一个连续函数分离（can be separated by a continuous function）。
\end{definition}

Urysohn 引理说明，如果 $X$ 中每一对无交的闭集能用无交的开集分离，那么它们也能用一个连续函数分离。其逆是显然的。因为若 $f\colon X\to[0,1]$ 是所述的连续函数，则 $f^{-1}\big([0,\frac{1}{2})\big)$ 和 $f^{-1}\big((\frac{1}{2},1]\big)$ 就是分别包含 $A$ 和 $B$ 的无交的开集。

现在可能会提出这样的问题：Urysohn 引理的证明是否能推广到正则空间上。既然在正则空间上能用无交的开集来分离点和闭集，那么是否也能用连续函数来分离点和闭集呢？

乍看起来，似乎可以像 Urysohn 引理的证明一样，先取一点 $a$ 和一个不包含这一点的闭集 $B$，如前面那样，定义 $U_{1}=X-B$。然后应用 $X$ 的正则性，选取包含点 $a$ 的一个开集 $U_{0}$，使得它的闭包含在 $U_{1}$ 中。但在紧接着下一步的证明中，就遇到了困难。设 $p$ 是序列中在 $0$ 和 $1$ 后面的一个有理数，要想找到一个开集 $U_{p}$，使得 $\overline{U}_{0}\subset U_{p}$，并且 $\overline{U}_{p}\subset U_{1}$。对于这一点，光靠正则性就不够了。

事实上，能够用一个连续函数来分离点和闭集，这比要求用无交的开集来分离它们的条件更强。我们把这一条件当作一个新的分离公理：

\begin{definition}\label{definition:ch04-completely-regular}
空间 $X$ 称为\textbf{完全正则的}（completely regular），如果每一个单点集是闭集，并且对于 $X$ 中的每一个点 $x_{0}$ 和不包含 $x_{0}$ 的任何一个闭集 $A$，存在一个连续函数 $f\colon X\to[0,1]$，使得 $f(x_{0})=1$ 和 $f(A)=\{0\}$。
\end{definition}

根据 Urysohn 引理，一个正规空间一定是完全正则的，并且一个完全正则的空间一定是正则的。这是因为，给定 $f$，集合 $f^{-1}\big([0,\frac{1}{2})\big)$ 和 $f^{-1}\big((\frac{1}{2},1]\big)$ 是分别包含 $A$ 和点 $x_{0}$ 的无交开集。于是，完全正则性就介于分离公理中的正则性和正规性之间。注意，定义中我们可以让函数 $f$ 将 $x_{0}$ 映射到 $0$，将 $A$ 映射到 $\{1\}$。因为函数 $g(x)=1-f(x)$ 就满足这一条件。但我们的定义要方便些。

在拓扑学的早年发展中，分离公理曾根据其性质由弱到强被排序为 $T_{1}$、$T_{2}$（Hausdorff）、$T_{3}$（正则性）、$T_{4}$（正规性）和 $T_{5}$（完全正则性）。字母``T''就是德语``Trennungsaxiom''，意思是``分离公理''。后来，当引入了完全正则的概念时，一些学者建议将其称为``$T_{3\frac{1}{2}}$ 公理''，因为它介于正则性和正规性之间。事实上，在一些文献中就用到这一术语。

与正规性不同，这个新的分离公理对于子空间和积空间而言有良好的表现。

\begin{theorem}\label{theorem:ch04-completely-regular-hereditary}
完全正则空间的子空间是完全正则的。完全正则空间的积空间是完全正则的。
\end{theorem}

\begin{proof}
设 $X$ 是完全正则的，$Y$ 是 $X$ 的一个子空间。令 $x_{0}$ 是 $Y$ 的一个点，$A$ 是 $Y$ 中的一个不包含点 $x_{0}$ 的闭集。那么 $A=\overline{A}\cap Y$，其中 $\overline{A}$ 表示 $A$ 在 $X$ 中的闭包。因此，$x_{0}\notin\overline{A}$。根据 $X$ 的完全正则性，我们可以选取连续函数 $f\colon X\to[0,1]$，使得 $f(x_{0})=1$ 和 $f(\overline{A})=\{0\}$。于是 $f$ 在 $Y$ 上的限制就是所要求的连续函数。

设 $X=\prod X_{\alpha}$ 是完全正则空间的一个积空间。令 $\boldsymbol{b}=(b_{\alpha})$ 是 $X$ 的一个点，$A$ 是 $X$ 中不包含点 $\boldsymbol{b}$ 的闭集。选取一个基元素 $\prod U_{\alpha}$，使得它包含点 $\boldsymbol{b}$ 并且与 $A$ 无交。于是除了有限多个 $\alpha$ 外，有 $U_{\alpha}=X_{\alpha}$。这有限多个 $\alpha$ 不妨设为 $\alpha=\alpha_{1},\alpha_{2},\cdots,\alpha_{n}$。给定 $i=1,\cdots,n$，选取连续函数
\begin{align*}
f_{i}\colon X_{\alpha_{i}}\longrightarrow[0,1],
\end{align*}
使得 $f_{i}(b_{\alpha_{i}})=1$ 及 $f_{i}(X-U_{\alpha_{i}})=\{0\}$。令 $\phi_{i}(\boldsymbol{x})=f_{i}(\pi_{\alpha_{i}}(\boldsymbol{x}))$，则 $\phi_{i}$ 连续地将 $X$ 映射到 $\mathbb{R}$，并且在 $\pi_{\alpha_{i}}^{-1}(U_{\alpha_{i}})$ 之外取值为零。于是积
\begin{align*}
f(\boldsymbol{x})=\phi_{1}(\boldsymbol{x})\cdot\phi_{2}(\boldsymbol{x})\cdot\cdots\cdot\phi_{n}(\boldsymbol{x})
\end{align*}
就是所求的 $X$ 上的连续函数，因为在 $\boldsymbol{b}$ 点处取值为 $1$，在 $\prod U_{\alpha}$ 外取值为零。
\end{proof}

\begin{example}\label{example:ch04-completely-regular-not-normal}
空间 $\mathbb{R}_{\ell}^{2}$ 和 $S_{\Omega}\times\overline{S}_{\Omega}$ 都是完全正则的，但不是正规的。因为它们都是完全正则空间（实际上是正规空间）的积空间。

要找一个正则却非完全正则的空间是比较困难的，已给出的大多数例子都很复杂，并要求熟悉基数理论。最近，Thomas [T] 构造了一个简单得多的例子，我们放在习题 11 中。
\end{example}

\subsection*{习题}

\begin{enumerate}
\item 验证 Urysohn 引理的证明，并且证明：对于给定的 $r$，
\begin{align*}
f^{-1}(r)=\bigcap_{p>r}U_{p}-\bigcup_{q<r}U_{q},
\end{align*}
其中 $p$ 和 $q$ 为有理数。
\item (a) 证明：多于一点的连通的正规空间是不可数的。

(b) 证明：多于一点的连通的正则空间是不可数的。[提示：任意可数空间都是 Lindelöf 空间。]
\item 在度量空间 $(X,d)$ 上令
\begin{align*}
f(x)=\frac{d(x,A)}{d(x,A)+d(x,B)},
\end{align*}
据此给出 Urysohn 引理的一个直接证明。
\item 我们曾经说过，若 $A$ 能表示成 $X$ 中可数个开集的交，则称 $A$ 是 $X$ 中的一个``$G_{\delta}$ 集''。

\textbf{定理}\quad 设 $X$ 是正规的。存在一个连续函数 $f\colon X\to[0,1]$，使得当 $x\in A$ 时，$f(x)=0$；当 $x\notin A$ 时，$f(x)>0$ 的充分必要条件是 $A$ 为 $X$ 中的一个闭的 $G_{\delta}$ 集。

满足此定理要求的函数称作是恰好在 $A$ 上\textbf{蜕化的}（vanish precisely on $A$）。
\item 证明：

\textbf{定理}（Urysohn 引理的强形式）\quad 设 $X$ 是一个正规空间。存在一个连续函数 $f\colon X\to[0,1]$，使得当 $x\in A$ 时 $f(x)=0$；当 $x\in B$ 时 $f(x)=1$；对于其他情形，有 $0<f(x)<1$ 当且仅当 $A$ 和 $B$ 是 $X$ 的两个无交的闭的 $G_{\delta}$ 集。
\item 一个空间 $X$ 被称为\textbf{完美正规的}（perfectly normal），如果 $X$ 的每一个闭子集都是 $X$ 的一个 $G_{\delta}$ 集。
\begin{enumerate}
\item 证明：每一个可度量化空间都是完美正规的.
\item 证明：完美正规空间都是完全正则的。因此，有时将完美正规条件称为``$T_{6}$ 公理''。[提示：设 $A$ 和 $B$ 是 $X$ 中的两个分离集。选取连续函数 $f,g\colon X\to[0,1]$ 使得它们分别恰在 $\overline{A}$ 和 $\overline{B}$ 上蜕化。考虑函数 $f-g$。]
\item 存在一个熟悉的空间，它是完全正则的，但不是完美正规的。这个空间是哪个呢？
\end{enumerate}
\item 证明每一个局部紧致的 Hausdorff 空间是完全正则的。
\item 设 $X$ 是完全正则的，$A$ 和 $B$ 是 $X$ 中的无交闭子集。证明：若 $A$ 是紧致的，则存在连续函数 $f\colon X\to[0,1]$ 使得 $f(A)=\{0\}$ 和 $f(B)=\{1\}$。
\item 证明 $\mathbb{R}^{J}$ 在箱拓扑下是完全正则的。[提示：指出只需要考虑箱邻域 $(-1,1)^{J}$ 与集合 $A$ 无交，点 $x$ 为坐标原点这样一个特殊情形。然后再应用一致拓扑下的连续函数也是箱拓扑下的连续函数这个结论。]
\item 证明以下结论：

\textbf{定理}\quad 每一个拓扑群都是完全正则的。

证明：在拓扑群 $G$ 中，设 $V_{0}$ 是单位元 $e$ 的一个邻域。一般地，选取 $e$ 的一个邻域 $V_{n}$，使得 $V_{n}\cdot V_{n}\subset V_{n-1}$。考虑所有二进制有理数 $p$ 构成的集合，即所有形如 $k/2^{n}$ 的有理数，其中 $k$ 和 $n$ 都是整数。对于 $(0,1]$ 中的每一个二进制有理数 $p$，给出一个开集 $U(p)$，其中 $U(p)$ 用以下方式归纳定义：$U(1)=V_{0}$，并且 $U\big(\frac{1}{2}\big)=V_{1}$。给定 $n$，假设对于 $0<k/2^{n}\leqslant 1$，$U(k/2^{n})$ 已有定义，然后对于 $0<k<2^{n}$ 定义
\begin{align*}
U(1/2^{n+1})=V_{n+1},
\end{align*}
\begin{align*}
U((2k+1)/2^{n+1})=V_{n+1}\cdot U(k/2^{n}).
\end{align*}
对于 $p\leqslant 0$，令 $U(p)=\varnothing$。对于 $p>1$，令 $U(p)=G$。证明对于所有的 $k$ 和 $n$，
\begin{align*}
V_{n}\cdot U(k/2^{n})\subset U((k+1)/2^{n}).
\end{align*}
后面的证法如 Urysohn 引理的证明。

这个习题来自 [M-Z]，读者可以在那里看到更多的关于拓扑群的结论。
\item 定义集合 $X$ 如下：对于每一个偶数 $m$，令 $L_{m}$ 表示平面中的线段 $m\times[-1,0]$。对于每一个奇数 $n$ 和整数 $k\geqslant 2$，设 $C_{n,k}$ 为平面中的线段 $(n+1-1/k)\times[-1,0]$ 和 $(n-1+1/k)\times[-1,0]$ 以及半圆
\begin{align*}
\{\,x\times y\mid(x-n)^{2}+y^{2}=(1-1/k)^{2}\quad\text{且 }y\geqslant 0\,\}
\end{align*}
的并。设 $p_{n,k}$ 是这个半圆的最高的点 $n\times(1-1/k)$。设 $X$ 是集合 $L_{m}$ 和 $C_{n,k}$，以及另外的两个点 $a,b$ 的并。通过取以下四种集合构成拓扑基，将 $X$ 拓扑化：
\begin{enumerate}
\item $X$ 与不包含点 $p_{n,k}$ 的水平开线段的交.
\item 从这些集合 $C_{n,k}$ 的某一个中删除有限多个点得到的集合.
\item 对于每一个偶数 $m$，$\{a\}$ 与 $X$ 中所有满足 $x<m$ 的点 $x\times y$ 构成的集合之并.
\item 对于每一个偶数 $m$，$\{b\}$ 与 $X$ 中所有满足 $x>m$ 的点 $x\times y$ 构成的集合之并.
\end{enumerate}
(a) 画出 $X$ 的草图。证明：这些集合构成了 $X$ 的一个拓扑基。

(b) 设 $f$ 是 $X$ 上的一个连续实值函数。证明：对于任意 $c$，集合 $f^{-1}(c)$ 是 $X$ 中的一个 $G_{\delta}$ 集。（这一点对于任意空间 $X$ 都是正确的。）证明由 $C_{n,k}$ 中满足 $f(p)\neq f(p_{n,k})$ 的点 $p$ 构成的集合 $S_{n,k}$ 是可数的。选取 $d\in[-1,0]$，使得直线 $y=d$ 与集合 $S_{n,k}$ 无交。证明：对于奇数 $n$，
\begin{align*}
f((n-1)\times d)=\lim_{k\to\infty}f(p_{n,k})=f((n+1)\times d).
\end{align*}
验证 $f(a)=f(b)$.

(c) 证明 $X$ 是正则的，但不是完全正则的.
\end{enumerate}